% Generated draft: requires mathematical and source review.
\documentclass[11pt]{article}
\usepackage[margin=25mm]{geometry}
\usepackage{fontspec}
\setmainfont{Latin Modern Roman}
\usepackage{amsmath,amssymb,mathtools,unicode-math}
\setmathfont{Latin Modern Math}
\usepackage{enumitem,needspace,xurl,hyperref}
\setlength{\parindent}{0pt}
\setlength{\parskip}{5pt}
\setlength{\emergencystretch}{3em}
\newcommand{\N}{\mathbb{N}}
\newcommand{\Z}{\mathbb{Z}}
\newcommand{\Q}{\mathbb{Q}}
\newcommand{\R}{\mathbb{R}}
\newcommand{\C}{\mathbb{C}}
\newcommand{\F}{\mathbb{F}}
\newcommand{\A}{\mathbb{A}}
\newcommand{\PP}{\mathbb{P}}
\newcommand{\E}{\mathbb{E}}
\newcommand{\eps}{\varepsilon}
\newcommand{\dd}{\,\mathrm{d}}
\begin{document}
\section*{UnsolvedMath ID: 30000424}
\textbf{Source problem code:} catalogue entry
\subsection{Definitions and mathematical statement}
Let \(k\) and \(c\) be the field and affine-space dimension intended in the source. The question asks whether every algebraic cone \(Z\subseteq k^c\) can be realized as
\[
 Z=V_R^*(M)
\]
for some module \(M\) over a complete intersection \(R\). The starred invariant \(V_R^*(M)\) is not defined in the supplied record and cannot be replaced by an unstarred support variety without confirmation.

In a standard affine convention an algebraic cone is the zero set \(Z(I)\subseteq k^c\) of a homogeneous ideal \(I\subseteq k[t_1,\ldots,t_c]\). The source does not specify algebraic closedness of \(k\), whether equality concerns \(k\)-rational zero sets or a scheme/geometric variety after scalar extension, or conventions about the empty cone and the origin. Set-theoretic realization does not automatically mean realization of a chosen nonreduced scheme structure.

In the usual local convention, a complete intersection is a Noetherian local ring with a presentation, possibly after completion, as a regular local ring modulo a regular sequence. To associate an affine space \(k^c\) of cohomological parameters, one ordinarily needs a residue-field identification with \(k\), codimension \(c\) and a choice of operators or coordinates. These data and any finite-generation condition on \(M\) are absent. One familiar support construction uses an annihilator of \(\operatorname{Ext}_R^*(M,k)\) under cohomological operators, but the record does not say that this is its starred invariant or prescribe such an action.

The printed existential wording can be read as
\[
 \forall Z\text{ an admissible cone in }k^c\ \exists(R,M),\quad Z=V_R^*(M).
\]
A stronger fixed-ring realization question would first fix every admissible complete intersection \(R\) and then ask \(\forall Z\ \exists M\). The source does not state whether \(R\) varies with \(Z\) or is fixed from omitted context. Both possibilities require the missing link between \(R,k,c\).

The literature assessment checked 2026-08-21 describes intermediate Ext-vanishing classes between complete intersections and Gorenstein rings, and an open classification direction. That does not explicitly assess the displayed cone-realization question or define \(V_R^*\). This mismatch is preserved rather than treated as proof that ordinary support-variety realization is open. The verification note says only that notation was delimited without changing the existential question. The open label remains dated curation metadata.
\subsection{Short English statement}
Can every algebraic cone be realized by the source's starred module invariant over a complete intersection? The invariant is undefined, and the field, codimension, module class and fixed-versus-variable ring quantifiers are missing. The literature assessment discusses a different Ext-vanishing classification direction, so it does not settle this formulation's status.
\subsection{Sources}
\begin{itemize}
\item Local UnsolvedMath registry, permanent ID 30000424, source code catalogue entry, ``Realization of Algebraic Cones over Complete Intersections,'' including original and cleaned realization question and the Ext-vanishing assessment.
\item \emph{Mini-Workshop: Thick Subcategories --- Classifications and Applications (2006)}, Oberwolfach Reports, pp.\ 461--510, DOI 10.4171/owr/2006/08: \url{https://doi.org/10.4171/owr/2006/08}. This is the sole supplied literature reference, labelled ``original OWR report.''
\item Literature triage checked 2026-08-21; local record updated 2026-09-16. No external source has been visited or independently checked here.
\end{itemize}
\end{document}
